
% !TeX spellcheck = pt_BR

\chapter*{Do que se trata a Dinâmica}
\addcontentsline{toc}{chapter}{Do que se trata a Dinâmica}

A Cinemática nos ensinou a descrever o movimento com precisão: como
dizer onde um corpo está, quão rápido vai, em que direção aponta sua
velocidade, como essa velocidade muda ao longo do tempo. É uma
linguagem poderosa --- e, como toda boa linguagem, deliberadamente
incompleta. Ela descreve \textit{o quê}, mas se recusa, por princípio,
a responder \textit{por quê.}

É exatamente essa pergunta que a Dinâmica se propõe a responder.

O conceito central é o de \textbf{força} --- e, com ele, uma ruptura
que levou quase dois mil anos para acontecer. Durante esse tempo, o
mundo acreditou em Aristóteles: a força seria a causa do movimento.
Para um corpo se mover, seria necessário que uma força agisse
continuamente sobre ele. Intuitivo e razoável, porém equivocado.

Foi Galileu quem percebeu a falácia, ao observar bolinhas deslizando
em planos inclinados perfeitamente lisos. Foi Newton quem a
formalizou, ao enunciar que a força não é a causa do movimento, mas
a causa da \textit{variação} do movimento --- da aceleração. Um corpo
em repouso permanece em repouso sem precisar de nenhuma força para
isso. Um corpo em movimento permanece em movimento, na mesma direção
e com a mesma velocidade, enquanto nenhuma força resultante agir sobre
ele. O que exige força é a mudança.

Esse deslocamento conceitual --- de força como causa do movimento para
força como causa da aceleração --- é a virada mais importante de todo
este livro, e ela está no coração do Segundo Ato.

O nosso estudo se dividirá em duas partes. Na primeira, construiremos
a \textbf{Dinâmica Vetorial}: as Leis de Newton, a força gravitacional,
a força de atrito, o impulso, o momento linear e suas leis de
conservação. É uma linguagem inteiramente vetorial --- forças têm
direção e sentido, e tratá-las como escalares seria perder precisão
onde a precisão mais importa. O coração dessa parte são as Leis de Newton: uma estrutura lógica que descreve, a partir de grandezas vetoriais, as causas do movimento.

Na segunda parte, introduziremos uma linguagem complementar e
inteiramente escalar: a \textbf{Dinâmica Escalar}, baseada nos
conceitos de trabalho e energia. Essa formulação não substitui as
Leis de Newton --- descreve a mesma realidade física por um caminho
diferente, mais conveniente em situações onde o que importa não é
o que acontece \textit{durante} o percurso, mas o que muda entre o
estado inicial e o estado final do sistema.

As duas linguagens juntas --- vetorial e escalar --- formam o arsenal
completo da mecânica clássica. Com elas, estaremos prontos para o
Terceiro Ato: aplicar tudo que foi construído ao mundo real, dos
planetas em órbita ao comportamento dos fluidos.